%% kapa: exam=1 round=30 year=2019 period=2 form=A subject=appraisal_law subject_ko=감정평가관계법규 seq=39 number=39
%% source: Q-Net 감정평가사 기출문제 (cst00309) artlSeq=5205205 file=extracted/1008/5205205/2019년 제30회 감정평가사 1차시험 문제지/2019년 제30회 감정평가사 1차시험 문제지/2019년도 제30회 감정평가사 1차 2교시 A형.hwp
%% conversion: kordoc 4.24.51 -> tools/gen1.py (auto); review=proofread-sonnet; flags=none
\begin{question}{39}
도시 및 주거환경정비법령상 정비구역에 관한 설명으로 옳지 않은 것은? (단, 조례는 고려하지 않음)
\begin{choices}
  \item 정비구역의 지정권자는 정비구역에서의 건축물의 최고 높이를 변경하는 경우에는 지방도시계획위원회의 심의를 거치지 아니할 수 있다.
  \item 정비구역의 지정권자는 정비사업의 효율적인 추진을 위하여 필요하다고 인정하는 경우에는 하나의 정비구역을 둘 이상의 정비구역으로 분할하는 방법으로 정비구역을 지정할 수 있다.
  \item 정비사업의 시행으로 토지등소유자에게 과도한 부담이 발생할 것으로 예상되는 경우 정비구역의 지정권자는 지방도시계획위원회의 심의를 거치지 아니하고 정비구역등을 해제할 수 있다.
  \item 주거환경개선사업은 사업시행자가 정비구역에서 정비기반시설 및 공동이용시설을 새로 설치하거나 확대하고 토지등소유자가 스스로 주택을 보전ㆍ정비하거나 개량하는 방법으로 할 수 있다.
  \item 정비구역등의 추진 상황으로 보아 지정 목적을 달성할 수 없다고 인정되어 정비구역등이 해제된 경우 정비계획으로 변경된 용도지역은 정비구역 지정 이전의 상태로 환원된 것으로 본다.
\end{choices}
\end{question}
\answer{3}
